% Fuente: Auxiliar 4, Pregunta 3.
{\color{MidnightBlue}
Agregamos una variable de holgura \(s_1\geq 0\) para la primera restricción
y variables de exceso \(s_2,s_3\geq 0\) para la segunda y tercera restricción:
\[
2x_1-10x_2+s_1=2,
\]
\[
2x_1+6x_2-s_2=24,
\]
\[
8x_1+4x_2-s_3=28.
\]

Por lo tanto, la forma estándar es:
\[
\begin{aligned}
\min \quad & 4x_1+15x_2\\
\text{s.a.}\quad
& 2x_1-10x_2+s_1=2,\\
& 2x_1+6x_2-s_2=24,\\
& 8x_1+4x_2-s_3=28,\\
& x_1,x_2,s_1,s_2,s_3\geq 0.
\end{aligned}
\]

Equivalentemente,
\[
x=
\begin{pmatrix}
x_1\\
x_2\\
s_1\\
s_2\\
s_3
\end{pmatrix},
\]
con
\[
A=
\begin{pmatrix}
2 & -10 & 1 & 0 & 0\\
2 & 6 & 0 & -1 & 0\\
8 & 4 & 0 & 0 & -1
\end{pmatrix},
\qquad
b=
\begin{pmatrix}
2\\
24\\
28
\end{pmatrix},
\qquad
c=
\begin{pmatrix}
4\\
15\\
0\\
0\\
0
\end{pmatrix}.
\]

Los vértices factibles son:

\[
A=\left(\frac{63}{8},\frac{11}{8}\right),
\]
obtenido de la intersección entre
\[
2x_1-10x_2=2,
\qquad
2x_1+6x_2=24.
\]

\[
B=\left(\frac{9}{5},\frac{17}{5}\right),
\]
obtenido de la intersección entre
\[
2x_1+6x_2=24,
\qquad
8x_1+4x_2=28.
\]

\[
C=(0,7),
\]
obtenido de la intersección entre
\[
x_1=0,
\qquad
8x_1+4x_2=28.
\]

Evaluamos la función objetivo
\[
z=4x_1+15x_2.
\]

\[
\begin{array}{c|c}
\text{Punto} & z=4x_1+15x_2 \\ \hline
A=\left(\frac{63}{8},\frac{11}{8}\right) 
& \frac{417}{8} \\[2mm]
B=\left(\frac{9}{5},\frac{17}{5}\right)
& \frac{291}{5} \\[2mm]
C=(0,7)
& 105
\end{array}
\]

El óptimo es
\[
(x_1^\star,x_2^\star)
=
\left(\frac{63}{8},\frac{11}{8}\right),
\]
con valor óptimo
\[
z^\star=\frac{417}{8}.
\]

La matriz \(A\) tiene \(3\) filas, por lo que cada base debe estar formada por
\(3\) columnas linealmente independientes.

En
\[
A=\left(\frac{63}{8},\frac{11}{8}\right),
\]
las restricciones activas son la primera y la segunda. Por lo tanto,
\[
s_1=0,
\qquad
s_2=0.
\]

Además,
\[
s_3=8x_1+4x_2-28.
\]

Entonces,
\[
s_3
=
8\left(\frac{63}{8}\right)
+
4\left(\frac{11}{8}\right)
-
28
=
\frac{81}{2}.
\]

Luego,
\[
x^A=
\begin{pmatrix}
x_1\\
x_2\\
s_1\\
s_2\\
s_3
\end{pmatrix}
=
\begin{pmatrix}
\frac{63}{8}\\
\frac{11}{8}\\
0\\
0\\
\frac{81}{2}
\end{pmatrix}.
\]

La base asociada es
\[
\mathcal{B}_A=\{x_1,x_2,s_3\}.
\]

En
\[
B=\left(\frac{9}{5},\frac{17}{5}\right),
\]
las restricciones activas son la segunda y la tercera. Por lo tanto,
\[
s_2=0,
\qquad
s_3=0.
\]

Además,
\[
s_1=2-2x_1+10x_2.
\]

Entonces,
\[
s_1
=
2
-
2\left(\frac{9}{5}\right)
+
10\left(\frac{17}{5}\right)
=
\frac{162}{5}.
\]

Luego,
\[
x^B=
\begin{pmatrix}
x_1\\
x_2\\
s_1\\
s_2\\
s_3
\end{pmatrix}
=
\begin{pmatrix}
\frac{9}{5}\\
\frac{17}{5}\\
\frac{162}{5}\\
0\\
0
\end{pmatrix}.
\]

La base asociada es
\[
\mathcal{B}_B=\{x_1,x_2,s_1\}.
\]

En
\[
C=(0,7),
\]
se activan las restricciones
\[
x_1=0,
\qquad
8x_1+4x_2=28.
\]

Por lo tanto,
\[
s_3=0.
\]

Además,
\[
s_1=2-2x_1+10x_2=2+70=72,
\]
\[
s_2=2x_1+6x_2-24=42-24=18.
\]

Luego,
\[
x^C=
\begin{pmatrix}
x_1\\
x_2\\
s_1\\
s_2\\
s_3
\end{pmatrix}
=
\begin{pmatrix}
0\\
7\\
72\\
18\\
0
\end{pmatrix}.
\]

La base asociada es
\[
\mathcal{B}_C=\{x_2,s_1,s_2\}.
\]
}

